\section{Descripción del problema}

Un proveedor de energía tiene un conjunto de centrales eléctricas, cada una con una capacidad diaria y un costo por MW producido, y debe planificar la producción de un horizonte de $n$ días. Cada cliente pide una cantidad de energía cada día y paga un precio propio por cada MW que recibe. El proveedor debe entregar a cada cliente al menos el $G$~\% de su demanda diaria y busca maximizar la ganancia neta: los pagos recibidos menos el costo de producción.

Además, para limitar el desgaste de las turbinas, una central produce en régimen alto cuando supera cierto porcentaje de su capacidad (el 80~\% en el enunciado), y no puede estar en régimen alto durante cierto número de días seguidos (2 en el enunciado). Tanto el porcentaje como el número de días pueden variar de una central a otra. El número de centrales, de clientes y de días es arbitrario.
